\chapter{CRONOGRAMA}

O cronograma a seguir apresenta uma previsão das principais etapas para o desenvolvimento do trabalho ao longo do período de execução do TCC. As atividades foram organizadas de forma a contemplar a revisão bibliográfica complementar, a construção do cenário de simulação, a realização dos experimentos, a análise dos resultados e a escrita final do trabalho.

\begin{quadro}[htbp]
\centering
\caption{Cronograma previsto para o desenvolvimento do TCC}
\label{qua:cronograma}
\small
\renewcommand{\arraystretch}{1.2}

\begin{tabular}{|p{9cm} | c | c | c | c | c | c|}
\hline
\textbf{Etapa} & \textbf{Jul} & \textbf{Ago} & \textbf{Set} & \textbf{Out} & \textbf{Nov} & \textbf{Dez} \\
\hline
Revisão bibliográfica complementar & X & X &  &  &  &  \\
\hline
Construção e ajuste da rede viária no \gls{sumo} & X & X &  &  &  &  \\
\hline
Definição dos perfis de motoristas &  & X & X &  &  &  \\
\hline
Definição dos fluxos e cenários de simulação &  & X & X &  &  &  \\
\hline
Execução das simulações &  &  & X & X &  &  \\
\hline
Coleta e organização dos dados &  &  &  & X & X &  \\
\hline
Análise dos resultados &  &  &  & X & X &  \\
\hline
Escrita e revisão do TCC &  & X & X & X & X & X \\
\hline
Preparação da apresentação e defesa &  &  &  &  &  & X \\
\hline
\end{tabular}

\par\vspace{0.25cm}
\begin{minipage}{\textwidth}
\footnotesize Fonte: Elaborado pelo autor.
\end{minipage}

\end{quadro}

A etapa de revisão bibliográfica complementar será voltada principalmente à busca de referências relacionadas à parametrização dos perfis de motoristas, às métricas de análise e ao uso do \gls{sumo} em estudos semelhantes. Embora a fundamentação teórica já tenha sido desenvolvida no PTC, essa revisão adicional será importante para apoiar a implementação prática da pesquisa.

Na etapa seguinte, será realizada a construção e o ajuste da rede viária no \gls{sumo}. Essa fase incluirá a obtenção da base inicial no OpenStreetMap, a conversão para o formato do simulador e os ajustes manuais necessários na rede, especialmente em sentidos de circulação, conexões e permissões de conversão.

Depois disso, serão definidos os perfis de motoristas, os níveis de fluxo e os cenários de simulação. Nessa fase, serão configurados os parâmetros de condução que representarão os perfis cauteloso, normal e agressivo, bem como os cenários de fluxo baixo, médio e alto.

Na sequência, serão executadas as simulações e realizada a coleta dos dados gerados pelo \gls{sumo}. Esses dados serão organizados para permitir a análise das métricas de desempenho, como tempo médio de viagem, velocidade média, tempo parado, formação de filas e vazão de veículos.

Por fim, será realizada a análise dos resultados, seguida da escrita e revisão final do TCC. A etapa final do cronograma será dedicada à preparação da apresentação e da defesa do trabalho.